\documentclass[11pt]{article}
\usepackage[margin=0.9in]{geometry}
\usepackage{amsmath,amssymb,amsthm,booktabs}
\usepackage{mathrsfs}
\usepackage[T1]{fontenc}
\usepackage{lmodern,microtype}
\usepackage[colorlinks=true,linkcolor=blue,urlcolor=blue]{hyperref}
\newtheorem{theorem}{Theorem}[section]
\newtheorem{lemma}[theorem]{Lemma}
\newtheorem{proposition}[theorem]{Proposition}
\newtheorem{corollary}[theorem]{Corollary}
\theoremstyle{remark}
\newtheorem{remark}[theorem]{Remark}
\newcommand{\norm}[1]{\left\|#1\right\|}
\title{Toward a four-dimensional Yang--Mills construction\\
\large Notebook 15: a uniform nonlinear constrained Wilson minimum}
\author{Research working document}
\date{16 October 2026}
\begin{document}
\maketitle
\begin{abstract}
We assemble a nonlinear small-field minimizer construction for the
bare Wilson action on homogeneous finite periodic regulators.
The constants and coarse-data analyticity radius are uniform in
the refinement and volume. The proof requires an equal-exponent
constrained Schauder estimate, an exact Wilson force factorization
with removable mesh factors, and a two-Hilbert-direction bound
for the nonlinear constraint curvature. These yield a strict
minimum in a fixed flat gauge slice for the corrected average.
A constraint-preserving gauge flow and a separate transfer to
ordinary group averages then give a minimum, modulo gauge,
among sufficiently small global physical C1-alpha potentials
with prescribed ordinary coarse average. This is a finite-regulator
classical ingredient for a constructive program. It does not yet
construct an interacting continuum quantum theory.
\end{abstract}

\noindent\textbf{Status and provenance.}
Research proofs subjected to separate internal mathematical audits;
no external certification or novelty claim. Source matching refers
to the deposited transcriptions at commit
\texttt{e392ef45013a581664bab76b90867170b30fe10e} of
\href{https://github.com/ethanwsussman/Yang--Mills-LLM-experiment/tree/e392ef45013a581664bab76b90867170b30fe10e}{Yang--Mills-LLM-experiment}.
Original CMP98 and CMP99b pages have not all been independently
recovered. The action and physical pairing match the deposited
four-dimensional bare Wilson conventions; their source-coordinate
Hessian is distinguished explicitly from the ordinary-potential
Hessian used in the nonlinear argument.

\section{The result and the remaining construction problem}
\label{nm15-overview}

Fix a matrix group $SU(N_c)$, normalized trace
$\tau=\operatorname{Tr}/N_c$, an odd block ratio $L\ge3$,
and $0<\alpha<1$. Fine spacing is $\eta=L^{-j}$ and the
retained coarse spacing is one. The regulators are compatible
finite periodic lattices. Potentials are Hermitian traceless
on the real locus and are complexified for analytic estimates.
The complete-field norm $N_\alpha$ controls the potential,
its physical first differences and their $\alpha$-H\"older
seminorms. Transposes always use physical sums. Fixed collar
enlargements and orientation multiplicities are included in
the constants, which may depend on $L,N_c,\alpha$.

The final theorem is Theorem~\ref{ot-minimum}. In words,
for sufficiently small arbitrary coarse data $B$ there is an
analytic small fine field with the exact ordinary homogeneous
constraint $M^{(j)}(U)=e^{iB}$, and its action is minimal among
all fields in a fixed small $N_\alpha$ ball satisfying that
constraint. Equality within the ball means gauge equivalence.
The corrected representative also obeys the sharper expansion
$A_*(B)=HB+O(\norm{B}_\infty^2)$ in $N_\alpha$, where $H$
is the flat minimizing extension. All these bounds are uniform
in the mesh and the compatible periodic volume.

There are three distinct operations in the proof. The flat
Green operator solves the nonlinear stationarity equation in
a fixed gauge slice. A genuine gauge flow covers nearby fields
with the \emph{same corrected constraint}. Finally, two explicit
gauge constructions relate the corrected and ordinary constraint
sets. Neither last step is a consequence of fixed-slice convexity
alone. The ordinary representative need not be in the same flat
gauge slice as the corrected one.

The flat operator identities, canceled Fourier aliases, and
physical exponential locality used below were derived in
Notebooks 8--12. The exact nonlinear constraint chart is from
Notebook 13. Notebook 14 supplies the ordinary/corrected gauge
relation and reference-coordinate comparison. This document
contains the additional estimates that close the nonlinear
minimum argument. Its symbols are restated at their use, but
the earlier flat alias estimates remain explicit prerequisites.

This closes a concrete classical small-field step. It leaves
open extension to local gauges and mixed retained histories,
control of large fields, integration of interacting fluctuations,
running-coupling estimates, and a nontrivial continuum observable
hierarchy. Euclidean invariance, positivity, and the physical
mass gap remain separate obligations. The gap in the constrained
flat auxiliary problem is not a physical Yang--Mills mass gap.

